% Standalone mathematical proof section for inclusion in the research report.
\section{An independent local proof of the 242-state construction}

Let $M$ be the inherited-label Combined243 machine. Its initial state is
$95$ and its separate halt target is $254$. Define $N$ by deleting state
$15$, changing $(11,0)$ from $(0,R,15)$ to $(0,R,27)$, and changing
$(27,1)$ from $(1,R,39)$ to $(1,R,22)$. All other retained rows agree.
The resulting machine has $242$ working states and $484$ binary rows.

\paragraph{The unused branch is proved, not inferred from a finite run.}
For a configuration $(q,h,T)$ of $M$, consider the invariant
\[
q\in\{18,47,65\}\implies T(h-1)=0,
\qquad q=27\implies T(h)=0.
\]
It holds vacuously at state $95$, with arbitrary tape and head. Every
incoming edge to any of $18,47,65$ writes zero and moves right, so it
establishes the first implication. The only incoming edges to $27$ are
$(18,0),(47,0),(65,0)$; each writes zero and moves left. The first
implication at its source therefore establishes the second implication
at the target. All other target states impose no condition. This is a
one-step inductive invariant of the literal original machine.

\paragraph{Five complete local blocks.}
Use the states other than $15,27$ as comparison checkpoints. Every
nonexceptional transition from a checkpoint is identical in both
machines and targets another checkpoint or the common halt. The
exceptional entry rows are exactly $(11,0),(18,0),(47,0),(65,0)$.

At $(11,h,T)$ with $T(h)=0$, split on $T(h+1)$. If it is zero, the
original state path is $11,15,21$, with head offsets $0,1,2$. The new
path is $11,27,32,45,21$, with head offsets $0,1,0,1,2$. All four reads
of the latter path are zero; its rows are
\[
(11,0)\mapsto(0,R,27),\quad
(27,0)\mapsto(0,L,32),\quad
(32,0)\mapsto(0,R,45),\quad
(45,0)\mapsto(0,R,21).
\]
Both endpoints are $(21,h+2,T)$. If $T(h+1)=1$, the paths are
$11,15,22$ and $11,27,22$; both endpoints are $(22,h+2,T)$ after two
steps. These two identities require no initial tape invariant because
the first transition itself writes the zero consulted by the detour.

For each $q\in\{18,47,65\}$, at a checkpoint with $T(h)=0$, the
invariant gives $T(h-1)=0$. The unchanged two-step path $q,27,32$
ends at $(32,h-2,T)$ in both machines. The altered row $(27,1)$ is not
used on these original entries. These are the remaining three identities.
Every stated block preserves the entire tape; only the displayed two
cells are inspected, so both exterior tape contexts remain arbitrary.

\paragraph{Halting equivalence.}
The initial configuration is a checkpoint and satisfies the original
invariant. At a paired checkpoint, take one identical ordinary step or
the appropriate complete local block. The resulting full configurations
again coincide. The invariant on the original side persists by its
one-step induction, so this comparison can be repeated indefinitely.
Every original block has one or two steps; every new block has one, two,
or four steps. The clocks are strictly positive and uniformly bounded,
and no exceptional block contains a halt state. Hence both executions
halt together or both have unbounded length. For every $h$ and binary
tape $T$,
\[
\operatorname{Halt}(M,95,h,T)
\quad\longleftrightarrow\quad
\operatorname{Halt}(N,95,h,T).
\]
This is an ordinary finite combinatorial proof with arithmetical
induction; it supplies a proof formalizable in PA and therefore ZFC.
No proof-assistant kernel replay is asserted here.

For a halting execution, let $Z$ count visits to $(11,0)$ whose next
cell is zero in the original execution. The exact time identity is
$T_N=T_M+2Z$, and in particular $T_M\leq T_N\leq2T_M$.
Combining the blank-input equivalence with the inherited calibration
of Combined243 gives the requested ZFC consistency equivalence for $N$.
The local reduction uses no register-frame or consistency assumption.

\paragraph{The reusable mechanism.}
The new state $27$ uses its old zero branch to consult a bit immediately
to the left, then restores the original head position before resuming.
Its original one branch was excluded by an inductive invariant and can
therefore implement the one branch of another decoder state. The
preceding bit is supplied by the incoming transition itself. Thus tape
context provides one part of the decoder control without consuming a
new working state. In this concrete table the same mechanism also
allows alternative 242-state machines deleting state $50$ or $171$;
these alternatives compete for the same formerly unused branch and
are not three independent simultaneous savings.
